\chapter{Ad-hoc Volumes and Streams}
\label{app:volumes}

Appendix~\ref{app:dagonfs} put the whole scratch directory of a workflow on a
DAGonFS mount. DAGonStar~3.0 goes further. It has one interface for ad-hoc
storage systems, so a workflow can start a volume, bind some of its tasks to
it, and stop it again. A volume can also be a node of the graph, created and
destroyed by tasks of the new type \texttt{TaskType.VOLUME}, and shared through
the ordinary \texttt{workflow://} schema. This appendix accompanies course
deck~A2. Its sources are \texttt{dagon/adhoc.py}, \texttt{dagon/volume.py} and
\texttt{dagon/streaming.py}, and the documents
\texttt{docs/adhoc\_filesystems.md} and \texttt{docs/volume\_tasks.md} of the
DAGonStar repository.

\section{One Interface for Many Storage Systems}

DAGonStar owns only the invocation, the readiness check, the shutdown and the
environment of the tasks bound to a resource. The storage software stays
separately installed: DAGonStar does not vendor, patch or deploy it. Five
providers are built in.

\begin{table}[h]
\centering
\begin{tabular}{@{}lll@{}}
\toprule
Provider & Access mode & Note \\
\midrule
\texttt{dagonfs} & posix, client or preload & frontends \texttt{fuse3}, \texttt{native}, \texttt{preload} \\
\texttt{beeond} & posix & only inside a Slurm allocation (Lecture~42) \\
\texttt{hercules} & preload & an interception library in \texttt{LD\_PRELOAD} \\
\texttt{capio} & preload & an interception library in \texttt{LD\_PRELOAD} \\
\texttt{dynostore} & object & keys in a catalog, no path to open \\
\bottomrule
\end{tabular}
\caption{The built-in ad-hoc providers of DAGonStar~3.0.0.}
\label{tab:adhoc-providers}
\end{table}

The access mode tells a task how to reach the data. A \emph{posix} resource is
an ordinary mounted directory. A \emph{preload} resource intercepts the file
calls of the programs its library is loaded into; the shell redirections of the
launcher are not intercepted, so a task writes through a program such as
\texttt{cp} or \texttt{tee}. A \emph{client} resource needs a client protocol,
and an \emph{object} resource is a store of keys.

\begin{sloppypar}
Every bound task receives
\texttt{DAGON\_ADHOC\_PATH}, \texttt{DAGON\_ADHOC\_ACCESS},
\texttt{DAGON\_ADHOC\_PROVIDER} and the same variables for each named resource,
\texttt{DAGON\_ADHOC\_<NAME>\_PATH}. With \texttt{export="GULF"} the path is also
exported under a name of your choice. Variables that look like credentials are
refused.
\end{sloppypar}

\begin{cautionbox}[Built-in names are protected]
In DAGonStar~3.0 the names \texttt{dagonfs}, \texttt{beeond},
\texttt{hercules}, \texttt{capio} and \texttt{dynostore} cannot be replaced,
not even with \texttt{register\_provider(..., replace=True)}. Register your own
adapter under another name.
\end{cautionbox}

\subsection{Scopes}

\begin{sloppypar}
A resource has one of three scopes. A \emph{workflow} resource, created with
\texttt{Workflow.create\_adhoc\_filesystem}, starts before the first task and
stops after the last; \texttt{bind(task)} selects its users, or
\texttt{default=True} binds every task. A \emph{task} resource, created with
\texttt{Task.create\_adhoc\_filesystem}, starts just before its task and stops
right after it, also when the task fails. A \emph{volume} resource is created by
a volume task (\Cref{sec:volume-tasks}).
\end{sloppypar}

\section{Writing an Adapter}

A new storage system needs one subclass of
\texttt{dagon.adhoc.AdHocFileSystem} and one call to
\texttt{register\_provider}. \Cref{lst:local-directory} serves a plain local
directory, so it runs on any laptop; the rest of this appendix uses it as a
real, if modest, provider. The same file binds a workflow-scope volume
\texttt{gulf} to two tasks and gives \texttt{mean\_CAP} a private task-scope
volume \texttt{tmp}.

\lstinputlisting[caption={An ad-hoc adapter and two scopes (\texttt{examples/A2-volumes/local\_directory.py}).},label={lst:local-directory}]{../../examples/A2-volumes/local_directory.py}

\lstinputlisting[style=bash,caption={Recorded output of \texttt{local\_directory.py}.}]{../../examples/A2-volumes/output/local_directory.txt}

The Capri mean, $(16.0+16.5+17.0+17.5)/4 = 16.75$, was read from the shared
volume, and both directories were removed when their resources stopped.

\section{Volume Tasks}
\label{sec:volume-tasks}

A volume task creates, or destroys, one volume as a node of the graph:

\begin{lstlisting}
gulf = DagonTask(TaskType.VOLUME, "Gulf", provider="local-directory", volume="gulf",
                 stage_in={"inputs/observations.csv": "workflow:///fetch/outputs/observations.csv"})
\end{lstlisting}

When it runs, it starts the volume, waits until the provider is ready and
copies each \texttt{stage\_in} entry into it. Each stage-in source creates a
dependency on its producer, as any \texttt{workflow://} reference does.

\begin{definitionbox}{A reference into a volume}
\sloppy
\texttt{workflow:///Gulf/gulf/means/CAP.txt} is the file \texttt{means/CAP.txt}
inside the volume \texttt{gulf} created by the volume task \texttt{Gulf}. The
schema keeps its meaning (workflow, task, path in the task's data space)
because the volume appears in \texttt{Gulf}'s folder under its own name. Only
the transfer changes: a command task is bound to the volume and its reference
is rewritten to the location inside it. Nothing is copied, so several tasks may
read and write the same volume.
\end{definitionbox}

A volume ends in one of two ways. A \emph{destroy} task
(\texttt{action="destroy"}) receives a dependency on every task that uses the
volume, copies its \texttt{stage\_out} entries into its own folder and releases
the volume; later tasks read the copies as ordinary files. Without a destroy
task the volume is released by \emph{reference counting}, as soon as the last
task that refers to it has finished, and in any case when the workflow ends,
also after a failure. Client and object volumes (DAGonFS \texttt{native},
DynoStore) have no path that a command could open, so references into them are
refused; such volumes are used through \texttt{stage\_in}, \texttt{stage\_out},
streams and \texttt{attach()}.

\section{Streams}

\begin{sloppypar}
A provider that supports streaming carries streams between tasks bound to the
volume. The producer publishes items to \texttt{\$DAGON\_STREAM\_OUT} with
\texttt{python -m dagon.adhoc stream publish}; the consumer names the volume in
the query of its reference, \texttt{workflow:///split/stations?stream=gulf},
and reads each item once with \texttt{stream items}. The consumer starts as
soon as the producer's launcher runs. The protocol is rename-free: the producer
writes a numbered marker after each complete item, so items become visible in
publication order, and the launcher writes an end marker with the command's
exit status, so a failed producer fails its consumers.
\end{sloppypar}

\section{The Station Means on a Volume}

\Cref{lst:volume-task} runs the station means of Lecture~12 on a volume. Three
tasks write into the same volume, a fourth streams the records of each station
to a fifth, and \texttt{Collect} stages the means out before releasing the
volume. A second run builds the same graph for the \texttt{dagonfs} provider
under portable emulation, without a destroy task.

\lstinputlisting[caption={Volume tasks and a stream (\texttt{examples/A2-volumes/volume\_task.py}).},label={lst:volume-task}]{../../examples/A2-volumes/volume_task.py}

\lstinputlisting[style=bash,caption={Recorded output of \texttt{volume\_task.py}.}]{../../examples/A2-volumes/output/volume_task.txt}

The means can be checked by hand from the dataset: Naples
$(14+15+16+17)/4 = 15.50$, Capri $16.75$, Ischia $(15+15.5+16+16.5)/4 = 15.75$.
The stream carried $3 \times 4 = 12$ records.

\begin{cautionbox}[Writes inside a volume make no edge]
A reference into the volume creates a dependency on the volume task, not on
the task that wrote the file. In the first run, \texttt{report} reads the
copies made by \texttt{Collect}, which waits for every user, so the order is
right. In the second run there is no destroy task: without the explicit edges
added by \texttt{add\_dependency\_to}, \texttt{report} could read the means
before they exist. This is the same rule as for BeeOND workflow volumes in
Lecture~42.
\end{cautionbox}

\section{Checkpoints, JSON, CWL and FAIR}

A volume does not survive the process. A create task's checkpoint is reused
only when it records the same specification and every user's checkpoint is
also complete; otherwise the volume is created again. A destroy task's
checkpoint is reused only after the SHA-256 digests of its staged outputs are
verified again. Both actions survive \texttt{as\_json()} and
\texttt{load\_json()}. CWL~v1.2 has no shared mutable storage, so the exporter
materializes the volume as a \texttt{Directory} output of the create step and
records the intent in a \texttt{dagon:Volume} hint; data written into the
volume by sibling steps is not exchanged through CWL. FAIR records mark the
edges served by a volume with \texttt{"transfer": "volume"}.

\begin{practicebox}
Under portable emulation a local directory stands in for the volume and no
provider is started. An emulated run checks the graph, the references and the
release; it is not evidence that the provider ran. The DAGonStar repository
records live runs of DAGonFS volume tasks in
\texttt{docs/performance/}.
\end{practicebox}
